评审会上，验收标准写了四条：模型要准；对两个群体要公平；要扛得住有人故意构造的输入；还不能泄露任何个人信息。

四条都合理，没有一条是刁难。可它们不是四个可以各自去优化的指标——它们之间横着几条硬约束，而那些约束不是``工程上难''，是\textbf{数学上不能同时满足}。两个群体的基率不同时，校准与等错误率不可兼得；差分隐私的每一次查询都在花掉一份不可再生的预算；某些分布上，标准准确率与鲁棒准确率之和有一个小于二的上界。这些都不是当前技术水平的局限，它们是可以写下来并证明的命题。

前面十六个部分讲的多是``怎样把一件事做好''。这一部分讲另一类问题：\textbf{当你想同时把几件事都做好，什么时候数学会直接告诉你不行，以及此时还剩下什么选择。}

\begin{center}
\begin{tikzpicture}[x=1cm,y=1cm,font=\small,>=stealth]
  \draw[black!50,fill=blue!9] (0,1.75) rectangle (2.9,3.05);
  \node[align=center] at (1.45,2.4) {同时要\\\(A\) 与 \(B\)};

  \foreach \i/\a/\b/\w in {
    0/隐私/效用/{预算随查询次数累加，用完为止},
    1/鲁棒/准确/{某些分布上两者之和有上界},
    2/校准/等错误率/{基率不同时不可兼得},
    3/不假设分布/强保证/{遗憾只对固定决策承诺},
    4/成对比较/绝对水平/{只定得出相差一个常数的尺度},
    5/均衡/最优/{无政府代价可以大于一}}{
    \pgfmathsetmacro{\yb}{4.32-\i*0.78}
    \draw[black!45,fill=black!4] (3.9,\yb) rectangle (12.9,\yb+0.6);
    \node[align=left,anchor=west,font=\footnotesize] at (4.05,\yb+0.3)
      {\a\ 与 \b\ ：\w};
    \draw[->,black!45] (2.95,2.4) -- (3.85,\yb+0.3);
  }
\end{tikzpicture}

\emph{六个单元是同一件事的六个实例。左边那句诉求每次都合理，右边那条结论每次都拦住它——拦住它的不是工程难度，是定理。}
\end{center}

\subsection{五条贯穿始终的判断}

\textbf{不可能与``还不够好''是两回事。}泛化界会随样本量收紧，隐私预算花光了还能重新分配，而公平性那三条判据在基率不同时的冲突不随任何东西变化——它是几个定义之间的代数矛盾，采样量根本不进入结论。读到一条限制时，第一件事是问它属于哪一类：数据不够、算力不够、还是定义本身在打架。前两类值得投入，第三类只能选择放弃哪一条。

\textbf{在被翻译成一个具体的量之前，``公平''``隐私''``鲁棒''都不是可以讨论的对象。}同一个模型可以在一个判据下公平、在另一个判据下不公平，而两边的数字都对——因为它们算的是方向相反的条件概率。隐私也一样：``去掉了姓名''与``某个人在与不在数据里输出分布差多少''是完全不同的承诺。这一部分每一章的第一步都是同一个动作：把那个词写成式子，然后才谈得上做到与否。

\textbf{保证的种类比保证的数值更要紧。}确定性保证说``条件满足就一定成立''；概率保证说``以 \(1-\delta\) 成立''，而随机性来自算法、不来自数据；经验证据说``我试了几种攻击都没打穿''。三者能支撑的动作差得很远，而最后一种根本不构成保证——攻击失败只是攻击失败。把它们混在同一份报告里而不加区分，是这个领域里最常见的失误。

\textbf{对手的存在改变问题的类型，不只是难度。}独立同分布这条假设一塌，``测试集上的期望误差''就失去意义，必须换成遗憾这种不对世界做任何假设的承诺。而一旦对手也在学习，收敛的对象就从最优点变成均衡，且收敛的往往只是时间平均——当前策略可以永远绕着均衡转圈。GAN 训练的不稳定不是实现问题，它是这个结构的直接后果。

\textbf{改不动结论，可以改问题。}均衡不理想就改规则而不是劝人改行为；三条公平判据不可兼得，就明确选定牺牲哪一条并把这个选择写进文档；隐私预算不够就减少查询，而不是偷偷调小噪声。这一部分给不出绕过定理的方法——那种方法不存在——它给的是把选择摆到台面上的语言。悄悄做的选择也是选择，只是没人为它负责。

\subsection{六章怎么安排}

\hyperref[ch:141]{``差分隐私：把``看不出是谁''写成一个会被花光的预算''}最适合先读，因为它的定义把这一部分的方法论演示了一遍：先说明``匿名化''为什么不是隐私，再把要求改写成``某一个人在与不在数据集里，输出分布差多少''。那个定义里的两个``任意''是它全部强度的来源，也正是第一部分讲量词时反复强调的事。三篇分别处理定义、噪声该加多少、以及预算怎样被消耗。

\hyperref[ch:142]{``对抗鲁棒性：邻域上的全称命题，训练只见过有限个点''}讲的是另一种``想同时要''。它最有用的一条是把 Lipschitz 常数变成可以计算的认证半径，从而把``试了几种攻击都没打穿''与``可以证明这个半径内不会翻转''分开——两者都能写进报告，能支撑的结论完全不同。

\hyperref[ch:143]{``公平性的不可能定理：几个都合理的要求，数学上互不相容''}是全部六章里定理最硬的一章，四行代数就把冲突证完。它的价值不在于告诉你该选哪个判据——那不是数学问题——而在于让``我们两个都要''这句话失去可能性，于是讨论被迫前进到``我们选哪一个、为什么''。

\hyperref[ch:144]{``在线学习与遗憾界：不假设数据怎么来，承诺才兑现得了''}与\hyperref[ch:146]{``博弈与均衡：均衡只保证没有人能单方面变好''}要连着读。前者建立``无遗憾''这个看似谦逊的承诺，后者用它回答一个大问题：当所有参与者都在学习时，长期会收敛到什么。答案是零和情形下时间平均收敛到极小极大均衡，而这条结论顺带给出了极小极大定理的一个构造性证明。博弈那一章还有两处与全书主线的接口：极小极大定理本质上是``什么时候可以交换 max 与 min''，Nash 均衡的存在性靠不动点定理——这两个动作在书末的收束部分都被单独讲过。

\hyperref[ch:145]{``从比较中学习：能定出的是差值，不是水平''}处理的是偏好数据。人只能可靠地说``A 比 B 好''，说不准``A 值 7.3 分''；而从比较里能恢复什么、恢复不出什么有明确答案。最要紧的一条是分数只在相差一个常数的意义下可辨识，所以任何依赖奖励绝对值的下游做法，都在使用一个未定义的量。

读这一部分时值得带着一个问题：手上这条``做不到''，究竟是谁做不到——是这个算法、这份数据，还是任何算法任何数据。三种答案对应三种完全不同的下一步。

\hyperref[ch:147]{``收束：每一章都懂了之后，它们合起来在回答什么''}紧接在后面，它把全书横竖各切一刀。这一部分正好为那里准备了最后一批材料：一个系统上线之后要面对的不只是分布漂移，还有拿着它的输出去调整自己行为的人。
